Since the self-consistent field has central symmetry, every electron state
has a definite value of its orbital angular momentum $l$.  For a fixed $l$,
the states of an individual electron are numbered in ascending order of energy
by the \emph{principal quantum number} $n$, whose possible values are
$n=l+1,l+2,\ldots$; this ordering convention is chosen to match the one used
for the hydrogen atom.  In a complex atom, however, the order in which levels
of different $l$ rise in energy generally differs from that in hydrogen.  For
hydrogen the energy is independent of $l$ altogether, and consequently states
with a larger $n$ always have a higher energy.  In a complex atom, by contrast,
the level with, for example, $n=5$, $l=0$ lies below the level with $n=4$,
$l=2$ (see \S~73 for further details).
